% ====================================================================
%  ABSTRACT
% ====================================================================
\chapter*{Abstract}
\addcontentsline{toc}{chapter}{Abstract}
\label{ch:abstract}

Spatial architectures for deep neural network (DNN) inference achieve
high throughput and energy efficiency by distributing computation across
large arrays of processing elements (PEs) with a multi-level on-chip
memory hierarchy.
The performance of these accelerators depends critically on the
\emph{mapping}---the schedule that determines how loop iterations are
tiled, ordered, and distributed across hardware resources.
Analytical frameworks such as Timeloop and FactorFlow enable fast
evaluation and automated exploration of candidate mappings, but they
operate on a \emph{single layer at a time}: every intermediate
activation tensor is written to off-chip DRAM after one layer completes
and read back when the next begins.
Since DRAM accesses are roughly two orders of magnitude more expensive
than on-chip accesses, this layer-by-layer data movement dominates the
energy budget of many workloads.

\emph{Layer fusion} addresses this bottleneck by executing multiple
consecutive layers in a single, interleaved schedule so that
intermediate activations remain entirely on-chip.
However, fusing $F$ layers multiplies the number of loop dimensions
from 6 to $6F$, introduces cascading tile-size dependencies between
producers and consumers, creates five distinct operand categories, and
causes the map space to grow to astronomical proportions---exceeding
$10^{463}$ candidates for an 11-layer fusion, over 400 orders of
magnitude larger than the single-layer case.

This thesis extends the FactorFlow framework to model, prune, and
evaluate fused-layer mappings.
Three modelling concepts are introduced: \emph{aliasing}, which
captures the dual role of each intermediate activation as both a
producer's output and a consumer's input; \emph{decoupling}, which
keeps per-layer weight dimensions independently tileable; and a
\emph{depth-first scheduling heuristic}, which enforces the
producer--consumer tile-size inequality at every hierarchy level,
pruning infeasible mappings before cost evaluation.
The single-layer cost model is generalised to a per-layer decomposition
of memory operations, energy (via Accelergy), and latency with
bandwidth-limited stall detection.

The extended framework is validated against the DepFiN accelerator, a
22\,nm depth-first CNN chip with a $16 \times 128$ PE array, executing
an 11-layer fused FSRCNN workload.
The model reproduces the compute-bound latency of 37.94\,M cycles
exactly, enforces zero intermediate DRAM traffic, and predicts
negligible on-chip stalls (24 cycles out of 37.94\,M, arising only
from boundary layers).

A systematic design-space exploration is then conducted across two
spatial architectures (DepFiN-like and Eyeriss-v1-like), four CNN
workloads spanning activation-dominant (FSRCNN, MC-CNN) and
weight-dominant (VGG16, ResNet18) profiles, and three execution modes
(full fusion, partial fusion, single-layer).
The results show that full fusion reduces DRAM traffic by 25--95\%
and delivers latency savings of up to 48\% for activation-dominant
workloads.
However, the enlarged on-chip memories required by fusion impose an
energy penalty of $1.3\times$ to $36.9\times$ relative to single-layer
execution.
When energy and latency are combined into the Energy-Delay Product
(EDP), full fusion wins in only 1 of 8 architecture--workload
combinations tested.
Partial (pairwise) fusion captures nearly all the latency benefit of
full fusion at a fraction of the memory cost, but the energy overhead
still favours single-layer execution in most cases.

The principal conclusion is that layer fusion is primarily a
\emph{latency optimisation}, not an energy optimisation: its benefit
grows with the activation-to-weight traffic ratio and diminishes as
on-chip memory overhead increases.
Partial fusion on a modestly sized architecture emerges as the most
practical deployment strategy, balancing throughput gains against
hardware cost.

\vspace{1em}
\noindent\textbf{Keywords:} deep neural networks, spatial architectures,
layer fusion, dataflow mapping, design-space exploration, FactorFlow,
DepFiN, analytical cost modelling.
